\documentclass{article}
\usepackage{amsmath,amssymb,hyperref,url}
\title{Corrected payments for learning with tail risk}
\author{}
\date{}
\begin{document}
\maketitle
\begin{abstract}
This account connects Q's payments for strategic resource allocation with P's sampled estimates of conditional value-at-risk.
The peers asked whether P's estimates could support Q's learning rule.
They found failures in Q's certificates and participation claim, then developed corrected payments, a fixed
risk approximation, and a replacement learning rule under added assumptions.
The corrected equilibrium balances payments exactly, while approximation errors bound losses in welfare,
incentives and participation under true risk preferences.
The thread remained a draft: this conditional result was supported, but broader convergence and feasibility
questions, and a full comparison with prior work, remained open when the research allowance ended.
\end{abstract}

\section{Introduction}
Q, by Garjani, Shokri and Kebriaei~\cite{Q}, studies how to share limited resources among people who act in their own interests.
Each person chooses an allocation and reports a price.
A payment rule is meant to make individual choices deliver the allocation that maximises total welfare.
Q claims that its quadratic payments achieve this at a unique Nash equilibrium, where nobody can gain by changing their own message alone.
It also claims balanced payments, voluntary participation, and convergence of a learning rule based on noisy observations.

P, by Wang, Huang, Xu and Johansson~\cite{P}, studies cooperating agents who care about unusually large losses.
Its risk measure is conditional value-at-risk (CVaR), which measures loss in a specified upper tail of the distribution.
Agents observe sampled loss values rather than risk values or gradients.
P uses empirical CVaR and random perturbations to estimate a gradient, then combines local updates with communication.
With fixed sample sizes and smoothing radii, its limiting objective is an expected empirical, smoothed approximation to population CVaR.

The proposed connexion was to put P's value estimates into Q's strategic learning scheme.
This required more than replacing an average by a risk measure.
The peers had to check whether the payments still aligned incentives, what objective the samples actually
represented, and whether the resulting learning rule converged.
The useful outcome was a narrower construction: implement a fixed risk approximation exactly, then bound its
economic error relative to true risk preferences.

\section{What was done}
Ada and Emmy first checked Q's algebra independently.
Both found that its strict matrix inequality conflicted with an equality needed for implementation.
They then divided the work: Emmy pursued an alternative implementation proof, while Ada studied the economic
effects of P's fixed sampling error.

Direct substitution into Q's printed payment revealed a further problem: an equilibrium participant could pay more than the value received.
Ada proposed a correction depending only on other participants' messages, and Emmy checked it.
Emmy then separated equilibrium implementation from learning convergence.
A weighted calculation supported the former, but a quadratic example defeated the automatic nonexpansiveness
needed for the original learning argument.
He therefore developed a different, nested proximal-point learner.

Meanwhile, Ada bounded the difference between the expected sampled objective and true CVaR directly in value,
without first estimating gradient error.
Emmy checked the bound and its consequences for welfare and incentives.
The peers also checked related literature and found existing work on sampled CVaR games and stochastic proximal methods.
Late in the thread, they resolved a local query problem: random perturbations could leave a participant's
feasible set, while shrinking the decision set could remove the zero allocation needed for participation.
Ada moved only the query points inward; Emmy checked the construction and its effect on curvature and error.

\section{Results}
\subsection{Corrected payments and implemented allocations}
Let \(N\ge2\) be the number of participants, indexed by \(n\).
Participant \(n\) chooses an allocation vector \(x_n\) in a compact convex set \(X_n\), and a nonnegative price vector \(p_n\).
The nonnegative vector \(c\) is total capacity, so feasible allocations satisfy \(\sum_n x_n\le c\), component by component.
Let \(\gamma>0\) be the payment scale and \(t_n^Q\) Q's printed quadratic payment, labelled \texttt{eq40} in its source.
A positive payment is a cost to the participant.
Ada's correction was
\[
 t_n^{\rm corr}=t_n^Q+
 \frac{\gamma N}{(N-1)^3}
 \left(\sum_{j\ne n}p_j\right)^\top
 \left(\sum_{j\ne n}x_j-\frac{N-1}{N}c\right).
\]
Here \(j\) indexes other participants, and the transpose denotes the ordinary inner product.
The added term cannot change anyone's best response because it does not depend on that person's own message.

The implementation proof needs additional assumptions.
Write \(V_n\) for a participant's valuation and \(f_n=-V_n\) for the corresponding cost.
The costs must have a common strong convexity bound \(m>0\), meaning a uniform positive quadratic lower curvature.
The proof also assumes valid rules for combining subgradients with local constraint conditions, and existence
of a planner Karush--Kuhn--Tucker (KKT) pair.
Such a pair consists of an optimal allocation and nonnegative capacity multipliers satisfying stationarity and complementary slackness.
A sufficient curvature condition is
\[
 m>\frac{\gamma(2N-1)^2}{4(N-1)^4}.
\]

Emmy multiplied the price components of the game's optimality operator by \((N-1)/N\).
This preserved individual optimality conditions and cancelled the interaction between common allocation and common price changes.
Completing squares in the remaining terms showed that two equilibria must have the same allocation and no price disagreement.
Conversely, a planner KKT pair supplied an equilibrium.
Thus every equilibrium implements the unique welfare allocation, although its common price need not be unique.
The full argument is in \path{emmy/research_note.md}, under ``A weighted operator repairs allocation
implementation''; Ada checked the calculation independently.

At these equilibria let \(p\) be the common price and \(\lambda=\gamma p\) the planner multiplier.
The corrected payment reduces to
\[
 t_n^{\rm corr}=\lambda^\top(x_n-c/N).
\]
Complementary slackness makes these payments sum to zero.
If \(0\in X_n\) and the concave valuation is normalised by \(V_n(0)=0\), planner stationarity gives
\[
 V_n(x_n)\ge\lambda^\top x_n,
 \qquad V_n(x_n)-t_n^{\rm corr}\ge\lambda^\top c/N\ge0.
\]
Participation therefore yields at least the normalised outside value.
These are equilibrium guarantees; balance away from equilibrium was not proved.

\subsection{From the sampled objective to true risk}
For participant \(n\), let \(J_n\) be the random loss and \(\beta_n\in(0,1)\) the upper-tail probability used for CVaR.
Write \(C_n(x)\) for population CVaR at allocation \(x\), and \(H_n(x)\) for the expected empirical CVaR from \(s_n\) samples.
Assume fresh independent, identically distributed samples from a distribution independent of the query point.
Assume also samplewise convex losses, bounded by \(|J_n|\le U_n\), with Lipschitz constant \(L_n\), which
bounds changes in loss per unit distance.
Define
\[
 e_n=\frac{U_n}{\beta_n}\sqrt{\frac{2\pi}{s_n}}.
\]
The threshold formula for CVaR makes \(H_n\le C_n\): minimising the sampled expression before taking expectation can only lower its value.
P's integrated distribution-function error bound supplies the other inequality, \(C_n-e_n\le H_n\).

Let \(\widetilde C_n\) be the average of \(H_n\) over a centred ball of radius \(\delta_n>0\), where those queries are allowed.
Convexity and the Lipschitz bound give
\[
 -e_n\le\widetilde C_n(x)-C_n(x)\le L_n\delta_n,
 \qquad w_n=e_n+L_n\delta_n.
\]
The number \(w_n\) bounds the variation of the approximation error between any two decisions.
Ada gave the proof in \path{ada/research-note.md}, under ``P's fixed-batch objective admits a sharper value
bound'' and ``Economic perturbation theorem''; Emmy independently checked them.

Use normalised true valuation \(V_n(x)=C_n(0)-C_n(x)\), and the same definition with \(\widetilde C_n\) for the surrogate valuation.
For satisfaction observations, the loss is negative satisfaction; deterministic payments enter additively by
CVaR's cash translation property.
At a corrected surrogate equilibrium, the true gain from any unilateral deviation is at most \(w_n\), and true
normalised utility is at least \(-w_n\).
Payments still balance exactly.
If \(\widetilde x\) is its allocation and \(x^*\) a true planner optimum, then
\[
 \sum_n C_n(\widetilde x_n)-\sum_n C_n(x_n^*)\le\sum_n w_n.
\]
The proof compares surrogate values at the two allocations; payment terms cancel in the corresponding unilateral comparisons.
If total population cost has strong convexity bound \(m_0>0\), the squared Euclidean allocation error is at most \(2\sum_n w_n/m_0\).
For P itself, including its extra comparison for the contracted domain, the value argument gives a limiting
error of order smoothing radius plus inverse square root of batch size.
It does not remove the inverse-radius variance cost of finite-time learning.

\subsection{Feasible local queries and a replacement learner}
Suppose a known ball of radius \(\rho_n>0\), centred at \(a_n\), lies inside \(X_n\) relative to its affine
hull, the smallest affine space containing the set.
For a direction \(u\) of norm at most one in that space, the peers used
\[
 q_n(x,u)=(1-\varepsilon_n)x+\varepsilon_na_n+\delta_nu,
 \qquad 0<\varepsilon_n<1,\quad 0<\delta_n<\varepsilon_n\rho_n.
\]
This is a convex combination of feasible points, so every query is locally feasible.
Decisions still range over the original set and include zero.
Let \(K_n\) be expected empirical CVaR averaged over these ball queries, and let \(R_n\) be the greatest
distance from \(a_n\) to a point in \(X_n\).
The economic bounds above then hold with
\[
 w_n^K=e_n+2L_n\varepsilon_nR_n+L_n\delta_n
\]
replacing \(w_n\).
If each sample loss has curvature \(m_n\), the new curvature is \(m_n(1-\varepsilon_n)^2\); its minimum must
exceed the implementation threshold.
Choosing \(\delta_n=\varepsilon_n\rho_n/2\) gives error of order \(s_n^{-1/2}+\varepsilon_n\).
Ada's continuation (E)--(H) contains the derivation, including the chain-rule factor in P's unbiased gradient estimator.
This construction assumes fresh losses can be queried at arbitrary locally feasible points.
It does not ensure that simultaneous experimental queries respect total capacity.

For fixed surrogate parameters, Emmy proved convergence of a changed learner.
It assumes a compact product strategy set with a price box containing a surrogate planner KKT equilibrium.
Each outer step solves a regularised version of the weighted equilibrium conditions approximately, using
projected stochastic gradient steps, P's value estimator, analytic payment derivatives and aggregate messages.
The regularisation makes each inner problem strongly monotone.

The expected squared inner error is bounded by a constant divided by the number of inner steps plus one.
At outer index \(k\), taking \(\lceil(k+1)^{2+\zeta}\rceil\) inner steps, for any fixed \(\zeta>0\), makes the
expected norm errors summable.
The proximal map's distance inequality then gives almost sure convergence; boundedness gives mean-square convergence.
The limiting allocation is the deterministic surrogate optimum, while the common price can have a random limit.
The complete recurrence proof is in Emmy's ``A conservative sampled-feedback convergence construction'', which Ada checked.
The record supplies no efficient total sample bound or stochastic performance benchmark.

\section{What did not hold}
Q's implementation equalities force zero curvature in a direction that changes all reported prices equally while leaving allocations fixed.
Its strict matrix inequality requires strictly negative curvature in every nonzero direction.
The peers' contradiction disproves simultaneous feasibility of those certificates, not existence of all possible implementing mechanisms.
The replacement argument does not validate Q's printed certificates.

The participation failure is concrete.
Ada took three participants, unit capacity and payment scale, allocations in \([0,1]\), and valuations
\[
 V_1(x)=10x-x^2/2,\qquad V_2(x)=V_3(x)=0.1x-x^2/2.
\]
Allocation \((1,0,0)\) with all prices equal to \(9\) satisfies the individual optimality conditions; positive
definite own-cost Hessians make them sufficient for equilibrium.
Yet Q charges the first participant \(10.5\) for value \(9.5\).
The corrected payments are \((6,-3,-3)\), which settle this failure.
Emmy checked a variant with coefficient \(1\) for the other valuations.
Both peers supplied exact rational checks in \path{ada/check_algebra.py} and \path{emmy/check_algebra.py}.

The same checks supported Emmy's separate quadratic example in which Q's proximal best-response map expands a
common-price perturbation for every positive proximal parameter.
Thus weighted monotonicity does not establish that map's Euclidean nonexpansiveness.
This rules out the proposed transfer of its certificate, without proving that every damped iteration diverges.
The nested learner supplies a different convergence result.

Q's curvature assumption concerns expected satisfaction; it does not supply the samplewise curvature needed here.
Nor is unique allocation the same as unique price.
Ada's strongly convex loss example produces a kink in CVaR at an interior optimum: the allocation is unique,
but every common price in \([1,3]\) is a valid multiplier.
Emmy checked it.
Fixed batches also retain a statistical bias, so convergence to their surrogate is not exact population-risk convergence.
No changing-parameter learning schedule was established for the corrected game.

\section{Why the thread stopped}
The verifier recorded a single decision, DRAFT.
Its reason recognised the connexion and accepted the corrected payments, weighted implementation, local
queries and replacement learner under their stated assumptions.
It also recognised that the account separated this conditional result from Q's invalid certificates and the
open questions about true-risk convergence, joint query feasibility and novelty.
It did not report a new algebraic defect.
The consolidated note says the research allowance ended without a final common readiness checkpoint covering every last entry.

The smallest recorded step towards restarting the publication question is to read the closest full papers and
compare assumptions and conclusions.
The ledger's search already found Wang, Shen and Zavlanos on empirical-CVaR game learning~\cite{wang22}, and
Cui and Shanbhag on stochastic proximal methods for hierarchical games~\cite{cui}.
Ada's note, incorporated by the ledger, also records work on learning Nash equilibria in risk-averse
games~\cite{wang24} and sample-average CVaR variational inequalities~\cite{saa}.
Her search checked landing pages and abstracts, not full proofs.
These references prevent novelty from resting merely on the estimator or the generic proximal method.
A focused comparison could establish whether the corrected transfers and explicit economic error bounds add a distinct result.
The record is too thin to settle that question; efficiency, original-rule convergence and jointly feasible
exploration require separate work.

\bibliographystyle{plain}
\bibliography{references}
\end{document}
